\documentclass[a4paper,11pt]{article} % Don't change this
\usepackage{fullpage} % Don't change this

\usepackage{titling} % Don't change this
\pretitle{\vspace{-6em}\begin{center}\larger[2]\bf} % Don't change this
\postauthor{\end{center}} % Don't change this
\preauthor{\vspace{-2em}\begin{center}\smaller[1]} % Don't change this
\postauthor{\end{center}\vspace{-5em}} % Don't change this

\usepackage[shortlabels]{enumitem} % Don't change this 
\setlist{itemsep=0em,topsep=0.3em} % Don't change this

%%% Some useful packages.
\usepackage{xcolor}
\usepackage{float}
\usepackage{amsmath,amssymb,mathtools}
\usepackage{graphicx}
\usepackage{subcaption}
\usepackage[hidelinks]{hyperref}
\usepackage{relsize}
\usepackage{cleveref}
\crefname{figure}{Fig.}{Figs.}

\usepackage[backend=bibtex,natbib=true,style=authoryear-icomp,maxnames=3,giveninits=true]{biblatex}
\citetrackerfalse
\bibliography{ref} % put bib entries in ref.bib

\usepackage[linesnumbered]{algorithm2e}
\usepackage{cleveref}
\crefname{equation}{Equation.}{Equations.}
\crefname{figure}{Figure}{Figures.}

\usepackage{lipsum}
\usepackage{wrapfig}
\usepackage{tikz}

%%% Add additional packages or macros below

\begin{document}

\title{STAT3007/7007 Project Report Template}
\author{Group X \\
Issac Newton (s12345678) $\cdot$
Albert Einstein (s23456781) $\cdot$
Alan Turing (s34567812) \\
John von Neumann (s45678123) $\cdot$ Kurt G\"odel (s56781234)}

\date{}
\maketitle

\begin{abstract}
	This is a \LaTeX\ template for your project report.
	Note the following: 
	(a) Change the title and group information above to your project title and
	group information.
	(b) Your report should not exceed \underline{12 pages} (excluding the references and
	appendices).
	You can change the section titles, but you should not alter the font size or
	the page size.
	It is perfectly fine to use less then 12 pages and conciseness is a marking
	criterion.
	(c) You may find it helpful to take a look at some NeurIPS papers at
	\url{https://papers.nips.cc/} and mimic their writing style.
	(d) You can do collaborative writing by uploading the template to \underline{Overleaf}:
	download \emph{STAT3007-7007\_Project\_Templates.zip},
	login to \url{https://www.overleaf.com},
	click \emph{Upload Project},
	then upload the zip file.
	Overleaf provides excellent guides on \LaTeX\ at
	\url{https://www.overleaf.com/learn}.
	Often, a quicker way to get the answer to a \LaTeX\ question is to ask ChatGPT;
	in many cases, it can give you a correct answer without much prompting.
	(e) When you submit your report, name it in the format of
	\emph{GX\_report\_AlanTuring.pdf}, where 
	\emph{GX} should be replaced by your group name, and \emph{AlanTuring} by the
	submitter's name.
	Only one group member needs to submit.
	(f) Have fun writing a nice report for your project findings.
\end{abstract}

\section{Introduction}

Describe your problem with a discussion on relevant background, motivation, and
significance in this section.

\section{Related Work}

Provide a discussion on related work in this section.
Use \verb+\citet+ to produce textual citations, and \verb+\citep+ to produce
parenthetical citations, as illustrated below:
\begin{itemize}
	\item \citet{mcculloch1943logical} proposed the linear threshold neuron as a
		mathematical model for biological neurons.
	\item One of the first mathematical models for biological neurons is the linear
		threshold neuron \citep{mcculloch1943logical}. 
\end{itemize}

\section{Methods}

Describe your methods in this section.
You should provide sufficient details so that your work is reproducible.

\section{Experiments}

Describe your experiment settings, datasets used, the results that you obtained,
and provide a discussion on the results. 
You should present your results in a concise manner.
For example, if you would like to include the learning curves for several
algorithms, generally you should present the curves in a single figure, rather
than producing one figure for one algorithm. 
Similarly, it is better to having a single table comparing the performance of
several algorithms, rather than having one table for each algorithm and
scattering them across your report.

\section{Conclusion}

Write your conclusion in this section.
You can discuss the limitations of your project here.

\section*{Contribution}

State the contribution of each group member.
Below is an example.
\begin{itemize}
	\item Newton: designed, implemented and evaluated method A; drafted content
		related to method A.
	\item Einstein: implemented and evaluated method B; drafted content related to
		method B.
	\item Turing: reviewed the literature and wrote Related Work.
	\item Neumann: collected and processed the datasets, set up the experiment
		environment; drafted content on these.
	\item G\"odel: compiled and polished the report.
\end{itemize}

\section*{Acknowledgement}

If anyone not in the group has collaborated with or advised you on the project, you
should fully acknowledge their contribution.

\section*{Appendix A.}
You can include any additional materials in one or more appendices.
Note that you should have all your key messages in the main text, and the
materials in an appendix should be supplementary but referenced in your main
text.
At most \underline{5 pages} are allowed for the appendices.

\smaller
\printbibliography

\newpage
Indicate whether all your group members consent for your report to be used as a
teaching resource by including one of the statement below.

\begin{center}
\emph{We give consent for this to be used as a teaching resource.}\\
\emph{We DO NOT give consent for this to be used as a teaching resource.}
\end{center}

\end{document}
